\chapter{DISCUSSÕES}
\label{cap:discussoes}

Este capítulo relaciona os resultados apresentados no
\autoref{cap:resultados} aos objetivos propostos e discute as
implicações e as limitações do sistema desenvolvido. A
\autoref{tab:cobertura_objetivos} relaciona cada objetivo específico
definido no \autoref{cap:objetivos} ao resultado observado no sistema
implementado.

\begin{table}[htb]
\centering
\caption{Cobertura dos objetivos específicos pelos resultados implementados}
\label{tab:cobertura_objetivos}
\begin{tabular}{p{0.04\textwidth}p{0.43\textwidth}p{0.45\textwidth}}
\hline
\textbf{Obj.} & \textbf{Objetivo específico} & \textbf{Resultado observado} \\
\hline
a & Protótipo ESP8266 com higrômetro, LEDs e \textit{buzzer} proporcional
ao risco &
Protótipo montado e em operação; LEDs acendem no nível correspondente;
\textit{buzzer} soa com frequência e cadência crescentes conforme a
gravidade do índice (80--84\%, 85--89\% e $\geq$90\%) \\
\hline
b & Integração com BNDMET/DECEA e OpenWeatherMap para precipitação
histórica e previsão &
Ambas as APIs integradas e em operação; acumulados de 24h, 7d e 30d via
BNDMET e previsão de 24h via OWM; proteção pré-NTP e \textit{retry}
acelerado implementados \\
\hline
c & Modelo de sete variáveis ponderadas com mecanismo de amplificação
$\times 1{,}20$ &
Modelo em operação; todos os sete componentes calculados e exibidos;
amplificação registrada no monitor serial e no \textit{dashboard};
protocolo de ruptura com \textit{cooldown} de três leituras validado \\
\hline
d & \textit{Backend} Node.js com banco de séries temporais e API
RESTful para recebimento, validação e persistência das leituras &
\textit{Backend} em operação com PostgreSQL e TimescaleDB; dados
recebidos, validados e persistidos; 41 campos por leitura;
autenticação JWT funcional \\
\hline
e & Interface \textit{web} com histórico, alertas, meteorologia e
gestão de usuários &
Seis telas implementadas e funcionais: \textit{dashboard}, Sensores,
Alertas, Relatórios, Usuários e Configurações; todos os fluxos de
sucesso e alternativos validados \\
\hline
f & Sistema de envio de alertas por e-mail acionado manualmente pelo
administrador, com formulário de composição, seleção de destinatários
por perfil e registro de auditoria &
Central de alertas operacional com formulário de composição, seleção
de destinatários por perfil, envio em massa e registro de auditoria
por disparo; e-mail HTML colorido por criticidade entregue aos
destinatários \\
\hline
\end{tabular}
\fonte{Elaborado pela autora.}
\end{table}

Como mostra a \autoref{tab:cobertura_objetivos}, todos os objetivos
específicos foram atendidos, e os resultados alcançados confirmam a
viabilidade técnica e econômica da solução proposta. O custo total do
protótipo de \textit{hardware} --- R\$~70,30 --- é expressivamente
inferior ao de instrumentos geotécnicos profissionais, como piezômetros
elétricos automatizados e inclinômetros de corda vibrante, cujos valores
de aquisição e instalação podem superar dezenas de milhares de reais por
ponto de monitoramento \cite{quispe2018instrumentacao}. Tecnologias de
baixo custo, quando bem integradas, são capazes de produzir fluxo
contínuo de dados geotécnicos e meteorológicos --- funcionalidade
ausente nos sistemas manuais que predominavam em Mariana e Brumadinho no
momento dos acidentes \cite{taliberti2025tragedia}.

A abordagem multidimensional do modelo de risco tem respaldo direto na
literatura. \citeonline{rico2008mine} identificaram que 39\% dos
rompimentos documentados resultaram da combinação de múltiplos fatores,
e não de uma causa isolada, o que justifica integrar variáveis
geotécnicas, pluviométricas históricas, prospectivas e atmosféricas em
um único índice. O peso maior atribuído à variável de lençol freático,
correspondente a 40\% do índice, foi definido pela autora por ser o
indicador mais direto da estabilidade do maciço, escolha coerente com o
nível d'água interno persistentemente elevado observado antes do colapso
de Brumadinho \cite{robertson2019feijaoson2019}. A inclusão da queda de
pressão atmosférica como indicador antecipatório representa uma
contribuição específica do modelo, fundamentada na relação entre queda
de pressão e aproximação de sistemas de mau tempo
\cite{noaa2024pressure, wmo2014guide}. Por se tratar de um sinal
indireto, recebeu o menor peso do modelo (5\%).

O índice de confiabilidade auditável diferencia o sistema da maioria das
soluções de baixo custo descritas na literatura, que emitem alertas sem
qualificar a incerteza dos dados subjacentes. Para qualquer leitura
histórica, o administrador pode consultar quais fontes estavam
indisponíveis e qual o impacto quantitativo de cada indisponibilidade
--- propriedade alinhada às recomendações do ICOLD sobre transparência
em sistemas de monitoramento \cite{icold2022tailings}. Além disso, a
decisão de processar o cálculo de risco integralmente no
microcontrolador garante que a análise continue ocorrendo localmente
mesmo diante de falhas de conectividade, o que é crítico para sistemas
destinados a ambientes industriais remotos.

Apesar desses resultados, sete limitações precisam ser registradas,
distribuídas entre o \textit{hardware}, o \textit{firmware}, o modelo de
risco e a operação da plataforma.

A primeira é espacial: o modelo atribui 40\% do índice de risco à
variável de lençol freático, estimada a partir de um único sensor em um
único ponto do maciço. Barragens de rejeitos reais têm extensão de
centenas de metros, e o nível de umidade varia de um ponto para outro da
estrutura --- regiões próximas a drenos internos ou a camadas de solo
com granulometria diferente podem estar em estados de saturação
completamente distintos do ponto monitorado. Uma leitura isolada não
captura esse comportamento distribuído, e é exatamente a variação ao
longo de toda a estrutura que determina se ela está estável ou em risco.
Um único ponto de medição, por mais representativo que seja na média,
não oferece essa visão completa \cite{machado2007}.

A segunda envolve o sensor e sua calibração. O higrômetro resistivo
FC-28 é um componente voltado a aplicações educacionais: a eletrólise
acelerada nos eletrodos em solo permanentemente úmido compromete sua
longevidade em instalações de longa duração. Agrava essa condição o fato
de a calibração ser feita manualmente via monitor serial, sem
procedimento periódico de verificação contra instrumento de referência e
sem registro da data da última calibração. Com o envelhecimento dos
eletrodos, a curva de resposta se desloca progressivamente, e o sistema
não possui mecanismo para detectar nem sinalizar esse desvio ao
administrador. Em uma implantação real, o instrumento seria cadastrado
no sistema com suas informações completas, incluindo os períodos de
aferição, manutenção, calibração e substituição previstos pelo
fabricante, e o sensor deveria ser substituído por instrumento com
certificação metrológica. Nesse cenário, o sistema estaria posicionado
como camada complementar à instrumentação geotécnica convencional
prevista pela Lei~n.º~14.066/2020 e pela NBR~13028:2017
\cite{brasil2020lei14066, abnt2017}.

A terceira é metodológica e incide sobre a equação do componente de
previsão pluviométrica. O fator~$F_{int}$, que compõe~$V_{ch.fut}$, é um
valor discreto em cinco classes (0,00 / 0,25 / 0,50 / 0,75 / 1,00), o
que introduz descontinuidades abruptas no índice: uma previsão de
4,9\,mm e outra de 0,1\,mm produzem o mesmo fator zero, enquanto o salto
de 4,9\,mm para 5,0\,mm causa um incremento instantâneo de 0,038 no
índice, independentemente das demais variáveis. Uma equação com
interpolação linear dentro de cada faixa, mantendo os valores extremos
atuais como âncoras, reduziria esse efeito sem alterar a estrutura
conceitual do modelo.

A quarta decorre do ciclo de vida do \textit{firmware}. Os \textit{buffers}
históricos --- o circular de umidade, usado em~$V_{taxa}$, e o de pressão
atmosférica, usado em~$V_{press}$ --- são inicializados com zeros a cada
reinicialização do dispositivo. Reinicializações involuntárias, por
estouro de \textit{heap} ou interrupção de energia, fazem com que os
primeiros ciclos pós-\textit{boot} calculem taxas de variação e quedas de
pressão comparando leituras reais contra zeros históricos, podendo gerar
valores artificialmente elevados para esses dois componentes até que os
\textit{buffers} completem um ciclo completo --- período que pode chegar a
três horas no caso do histórico de pressão.

A quinta diz respeito ao escopo de validação do modelo. Os pesos das
sete variáveis foram arbitrados pela autora a partir da abordagem
multidimensional de \citeonline{rico2008mine} e não passaram por processo
de calibração estatística sobre casos históricos de ruptura em barragens
de rejeitos brasileiras. Quanto aos limiares, o de 24 horas apoia-se em
\citeonline{rico2008mine} e no Sistema Alerta Rio; o de 30 dias foi
definido a partir do histórico da estação D6594; e os limiares de 25\%
e 30\% de umidade, de 7 dias (150\,mm) e as faixas de intensidade da
previsão foram adotados pela autora, sem ajuste às características
específicas dos rejeitos de minério de ferro predominantes no contexto
de referência. A representatividade da estação D6594 para um ponto
específico de monitoramento, a localização ideal da sonda no perfil do
maciço e a calibração dos limiares para diferentes tipos de rejeito
demandariam validação geotécnica especializada antes de qualquer
aplicação operacional.

A sexta afeta a operação cotidiana da plataforma. Quando o dispositivo
perde conectividade, os ciclos de leitura são calculados localmente, mas
não persistidos para envio posterior, o que gera lacunas na série
temporal registrada no banco. Além disso, o sistema não emite
notificação automática ao administrador quando o nível de alerta evolui
de AMARELO para VERMELHO: o sino no \textit{header} exibe o
\textit{badge} apenas para quem estiver com o painel aberto no momento
da atualização periódica. Em um cenário operacional real, em que o
responsável pelo monitoramento não permanece continuamente diante da
interface, isso representa um risco de resposta tardia. Conectividade
redundante --- Wi-Fi associada a módulo 4G --- e mecanismos de alerta
proativo por \textit{push} ou SMS mitigariam essas duas fragilidades.

Por fim, a sétima envolve a ausência de monitoramento de consumo
energético: o protótipo não incorpora sensor de corrente ou tensão, o
que impede alertas de bateria baixa em cenários de alimentação autônoma
(por exemplo, painel solar), aplicáveis fora do escopo de teste em
bancada aqui validado.